\begin{frame}{Takeaways — Day 1}

  \setlength{\itemsep}{2.5pt}
  \setlength{\parsep}{0pt}

  \begin{enumerate}
    \item \textbf{Defining an information standard}
    \begin{itemize}
      \item interoperability is a community quality; a standard is what they decide to make \emph{testable}
    \end{itemize}
    \pause
    \item \textbf{The best practice stack}
    \begin{itemize}
      \item L0 concepts (SKOS) $\rightarrow$ L1 open-world vocabulary $\rightarrow$ L2 profile $\rightarrow$ L3 \emph{derived} artifacts; schema $\neq$ meaning
      \item MOF: choose a pivot model, don't chase one; validate from the profile (SHACL)
      \item proven in practice: Netbeheer Nederland runs this with CIM/CGMES
    \end{itemize}
    \pause
    \item \textbf{openADR 3.0 through that lens}
    \begin{itemize}
      \item the normative core of the standard is a \emph{serialization}: the OpenAPI YAML,
            with the model defined inside it
      \item the three requirements of the day attach to exactly that choice:
            multiple use cases, technology-independence, no empty sets
    \end{itemize}
  \end{enumerate}
\end{frame}
